% !TEX TS-program = latex
% !TEX TS-program = xelatex
% -*- mode : latex; coding: utf-8 -*-

%% 재단 맡길 때에는 showtrims을 켜서 a4로 출력 후 잘라내면 됨.
% \documentclass[11pt,a4paper,oneside,openany,reqno,chapter,nokorean,showtrims]{oblivoir}
\documentclass[11pt,a4paper,oneside, openany, reqno, chapter, nokorean]{oblivoir}
\setstocksize{297mm}{210mm}
% \settrimmedsize{\stockheight}{\stockwidth}{*}

%\usepackage[utf8]{inputenc}
\usepackage[dbl4x6]{fapapersize}
\usepackage{fapapersize}
\usefapapersize{210mm,297mm,40mm,30mm,30mm,32mm}

%% figure setting
% \usepackage{geometry}

\usepackage{textcomp}

\usepackage{pgfplots}
\pgfplotsset{width=7cm, compat=1.16}


\let\interHchar\undefined
\let\interhword\undefined
\usepackage[default]{dhucs-interword}
% \usepackage{dhucs-gremph}
\usepackage{dhucs-enumerate}
\usepackage{ifpdf}
\usepackage{url}
\usepackage{rotating}
\usepackage{titletoc}
\usepackage{titlesec}
% \usepackage[toc,page,header]{appendix}
% center the toc heading
% \titleformat{\chapter}[display] {\normalfont\bfseries}{}{0pt}{\Huge}
\titleformat{\chapter}[display]
  {\normalfont\bfseries}{}{0pt}{\Huge}

%\usepackage{cleveref}

% Define formats

\usepackage{lipsum} 

\newcommand{\BibTeX}{\textsc{Bib}\TeX}

\chapterstyle{default}
\usepackage{titletoc}% http://ctan.org/pkg/tocloft
\def\cftchapteraftersnum{.}
\renewcommand{\cftchapterdotsep}{\cftdotsep}% Chapters should use dots in ToC

\makeatletter
\def\@makechapterhead#1{%
  %%%%\vspace*{50\p@}% %%% removed!
  {\parindent \z@ \raggedright \normalfont
    \ifnum \c@secnumdepth >\m@ne
        % \huge\bfseries \@chapapp\space \thechapter
        \par\nobreak
        \vskip 20\p@
    \fi
    \interlinepenalty\@M
    \LARGE \bfseries \thechapter. #1\par\nobreak
    \vskip 20\p@
  }}
\def\@makeschapterhead#1{%
  %%%%%\vspace*{50\p@}% %%% removed!
  {\parindent \z@ \raggedright
    \normalfont
    \interlinepenalty\@M
    \LARGE \bfseries  #1\par\nobreak
    \vskip 40\p@
  }}
\makeatother

\ifpdf

\hypersetup{
colorlinks = true
} 
\fi

\usepackage{amsmath,amsthm,amssymb,amsxtra}
\usepackage{graphicx}
\newsubfloat{figure}
\usepackage{epsfig}
\let\footruleskip\undefined
\usepackage{fancyhdr}
\usepackage{environ}
\usepackage{siunitx}
\ifpdf
% enable hyper links
\hypersetup{
 pdftitle={Thesis},
 pdfauthor={Author name},
 pdfkeywords={j/psi,charmed meson},
 citecolor=black,
 %citecolor=darkblue,
 filecolor=black,
 linkcolor=black,
 %inkcolor=darkblue,
 urlcolor=blue,
 colorlinks=true}
\fi
\usepackage{seoultech} % 서울과학기술대학교 양식

\setlength\parindent{2pc} % 들여쓰기 1글자

\newtheorem{theorem}{Theorem}[chapter] % Theorem 1.1, Theorem 3.4
\newtheorem{corollary}[theorem]{Corollary} % Corollary는 독자번호 없음.
\newtheorem{lemma}[theorem]{Lemma}
\newtheorem{proposition}[theorem]{Proposition}
\newtheorem{question}{Question} % 독자적인 번호를 갖는다. Question 1, Question 2
\newtheorem{step}{Step}

% 번호를 붙이지 않는 `정리' 만들기
\newtheorem{lemmawon}{Lemma}
\newtheorem{claim}{Claim}

% 특별한 이름을 갖는, 번호가 붙지 않는 Theorem 만들기
\newtheorem{questionA}{Question~A}
\newtheorem{mainthmA}{Main Theorem~A}

% 본문의 글씨가 기울어지지 않는 환경
\theoremstyle{definition}
\newtheorem{definition}[theorem]{Definition}

% 화면에 나오는 이름이 이탤릭체인 환경
\theoremstyle{remark}
\newtheorem{remark}[theorem]{Remark}
\newtheorem{remarkof}{Remark}[theorem] % 어떤 Theorem의 아래에 Remark
\newtheorem{remarkwon}{Remark}

% 수식 번호 붙이기 (1.1), (3.3)
\numberwithin{equation}{chapter} 

%책갈피용
%\usepackage[bookmarks=true]{hyperref}
\usepackage{bookmark}


\setcounter{secnumdepth}{4}
\setcounter{tocdepth}{2}


%표 정렬용
\usepackage{multirow}
\renewcommand{\multirowsetup}{\centering}

%그림위 글자 기능(인준서 도장용)
\usepackage[abs]{overpic}
\setlength\unitlength{1mm}

\usepackage{setspace}
%\doublespacing
\usepackage{array}
\usepackage{makecell}
\usepackage{multirow}
\usepackage{parallel,enumitem}
\usepackage{graphicx}
\usepackage{tikz}
\usetikzlibrary{shapes,decorations.pathmorphing}
\usepackage{amssymb}% http://ctan.org/pkg/amssymb
\usepackage{booktabs}
% \usepackage{adjustbox}
\usepackage{amsmath}
\usepackage{algorithm}
\usepackage{algorithmicx}
\usepackage[noend]{algpseudocode}
\usepackage[flushleft]{threeparttable}

\usepackage{subcaption}
\usepackage{siunitx}
\usepackage{tabularx}
\usepackage{enumitem}

\usepackage{amsthm}
\usepackage{stmaryrd}
\usepackage{amsfonts}
\usepackage{xcolor}

\usepackage{array}
\newcolumntype{L}[1]{>{\raggedright\let\newline\\\arraybackslash\hspace{0pt}}m{#1}}
\newcolumntype{C}[1]{>{\centering\let\newline\\\arraybackslash\hspace{0pt}}m{#1}}
\newcolumntype{R}[1]{>{\raggedleft\let\newline\\\arraybackslash\hspace{0pt}}m{#1}}

\usepackage{fontspec}
\newfontfamily\unbatang{UnBatang.ttf}
\usepackage{kotex}

% \setmainfont{Times New Roman}
% \setsansfont{Arial}
% \setmonofont{Courier New}

%%%%%%%%%%%%%%%%%%%%%%%%%%%%%%%%%%%%%%%%%%%%%%%%%%%%%%%%%%%%%%%%%%%%%%%%%%%
% to be able to draw some self-contained figs
\usepackage{tikz}
\usepackage{amsmath}
%%%%%%%%%%%%%%%%%%%%%%%%%%%%%%%%
% inlined bib file

%%%%%%%%%%%%%%%%%%%%%%%%%%%%%%%%%%%%%
\newcommand{\minus}{\scalebox{0.75}[1.0]{$-$}}
\theoremstyle{definition}
\theoremstyle{remark}
\theoremstyle{definition}

\algnewcommand\algorithmicinput{\textbf{Input:}}
\algnewcommand\INPUT{\item[\algorithmicinput]}
\algnewcommand\algorithmicoutput{\textbf{Output:}}
\algnewcommand\OUTPUT{\item[\algorithmicoutput]}

%%%%%%%%%%%%%%%%%%%%%%%%%%%%%%%%%%%%%%%%%%%%%%%%%%%%%%%%%%%%%%%%%%%%%%%%%%%%%%



\DeclareMathOperator*{\argmax}{argmax} % thin space, limits underneath in displays
  
\begin{document}
  %% 정보
  \EnglishTitle{Feature-Aware Numerical Representation for Retrieval-Augmented Medical Tabular Classification (FiTabr)}
  %% 본문이 영문일 경우, 국문 제목은 빈칸으로 둘 것
  \KoreanTitle{}
  \AuthorKoreanName{저자 한국 이름}
  \AuthorEnglishName{Kurbanov Elvin Vagibovich}
  \KoreanNameofDegree{공학 석사}
  
  \EnglishNameofDegree{Master of Science}
  \EnglishNameofPaper{Master Thesis}

  \GraduateDate{2026년 2월}
  \GraduateDateEnglish{February 2026}
  \EnglishDepartmentName{Department of Computer Science and Engineering}
  \KoreanDepartmentName{학과 한글 이름}
  \RefereeChief{심사위원장 이름}
  \RefereeSecond{심사위원1 이름}
  \RefereeThird{심사위원2 이름}
  \AdvisorKoreanName{지도교수 한글 이름}
  \AdvisorEnglishName{Taesoo Kwon}

%% 앞표지
\makefrontcover


%% 
\englishproposal

%% 국문 제출 승인서. 국문 인준서의 경우 업데이트가 되지 않아서 추가 작업 필요.
%\makeapproval

%% 영문 제출 승인서
\makeenglishapproval

\pagenumbering{gobble}


\cedp
  
 %% 본문 이전의 쪽들
\frontmatter 


\addcontentsline{toc}{chapter}{Summary}
\markboth{Summary}{}
\begin{EnglishSummary}

With the increasing application of artificial intelligence in healthcare, medical tabular data have become an important source for clinical prediction tasks. Such datasets typically contain heterogeneous patient-related variables, including demographic characteristics, laboratory measurements, and physiological indicators. However, learning from medical tabular data remains challenging because of limited sample sizes, heterogeneous feature types, class imbalance, and variation in the predictive relevance of individual variables. Recently, retrieval-augmented tabular learning methods such as TabR have demonstrated promising performance by incorporating information from similar training samples into the prediction process. However, the numerical representation used before retrieval does not explicitly model sample-specific variation in feature emphasis. Therefore, this thesis investigates feature-aware numerical representation for retrieval-augmented medical tabular classification. A feature-aware extension of TabR, termed FiTabR, is proposed by introducing a lightweight attention module before the TabR encoder. The proposed module generates sample-specific attention weights for numerical features and applies residual feature modulation while preserving the original retrieval, contextual aggregation, and prediction mechanisms of TabR. In addition, ablation studies are conducted to examine the contribution of residual preservation, attention scaling, and learned feature weighting to the behavior of the proposed architecture. The learned attention distributions are further analyzed using their mean and standard deviation to characterize feature emphasis and its correspondence with established medical domain knowledge. Experimental results across multiple medical tabular datasets show that FiTabR improves TabR under certain dataset conditions, while the magnitude of improvement varies across datasets. The ablation and attention analyses further indicate that the contribution of the proposed architectural components and the learned feature emphasis is dataset dependent. These findings demonstrate the potential of feature-aware numerical representation for retrieval-augmented tabular learning while also identifying its limitations across different medical data characteristics.


\end{EnglishSummary}

    
\pagenumbering{gobble}
 
    

\cedp
  
  %% 본문 이전의 쪽들
\frontmatter 

  %% TOC/ LOF / LOT 끝의 * 지우면 안됨. ( 목차에서 명단 삭제 필요 )

\pagestyle{plain}%
	%% 목차
\clearpage
\phantomsection 
\renewcommand{\contentsname}{Table of Contents}
\setcounter{tocdepth}{1}
\tableofcontents*
    %% 표목차

\clearpage
\phantomsection
\listoftables*
\addcontentsline{toc}{chapter}{List of Tables}

\clearpage
\phantomsection
\listoffigures*
\addcontentsline{toc}{chapter}{List of Figures}
        
  %%% 본문 시작
  \mainmatter
  \linespread{2}
  %\SetHangulspace{2.2}{2.0}

  %% 쪽번호가 아래 중앙에 붙는다.
  % \pagestyle{fancy}%
  % \lhead{\leftmark}%
  % \rhead{}%
  % \chead{}%
  \cfoot{\thepage}%
  \renewcommand{\headrulewidth}{0pt}

    \chapter{Introduction}
    This document is a model and instructions for \LaTeX.
    Please observe the conference page limits. 
    
    \chapter{Ease of Use}
    
    \section{Maintaining the Integrity of the Specifications}
    
    The class file is used to format your paper and style the text. All margins, column widths, line spaces, and text fonts are prescribed; please do not alter them. You may note peculiarities. For example, the head margin measures proportionately more than is customary. This measurement and others are deliberate, using specifications that anticipate your paper as one part of the entire proceedings, and not as an independent document. Please do not revise any of the current designations.
    
    
    \chapter{Prepare Your Paper Before Styling}
    Before you begin to format your paper, first write and save the content as a 
    separate text file. Complete all content and organizational editing before 
    formatting. Please note sections \ref{AA}--\ref{SCM} below for more information on 
    proofreading, spelling and grammar.
    
    Keep your text and graphic files separate until after the text has been 
    formatted and styled. Do not number text heads---{\LaTeX} will do that 
    for you.
    
    \section{Abbreviations and Acronyms}\label{AA}
    Define abbreviations and acronyms the first time they are used in the text, 
    even after they have been defined in the abstract. Abbreviations such as 
    IEEE, SI, MKS, CGS, ac, dc, and rms do not have to be defined. Do not use 
    abbreviations in the title or heads unless they are unavoidable.
    
    \section{Units}
    \begin{itemize}
    \item Use either SI (MKS) or CGS as primary units. (SI units are encouraged.) English units may be used as secondary units (in parentheses). An exception would be the use of English units as identifiers in trade, such as ``3.5-inch disk drive''.
    \item Avoid combining SI and CGS units, such as current in amperes and magnetic field in oersteds. This often leads to confusion because equations do not balance dimensionally. If you must use mixed units, clearly state the units for each quantity that you use in an equation.
    \item Do not mix complete spellings and abbreviations of units: ``Wb/m\textsuperscript{2}'' or ``webers per square meter'', not ``webers/m\textsuperscript{2}''. Spell out units when they appear in text: ``. . . a few henries'', not ``. . . a few H''.
    \item Use a zero before decimal points: ``0.25'', not ``.25''. Use ``cm\textsuperscript{3}'', not ``cc''.)
    \end{itemize}
    
    \section{Equations}
    Number equations consecutively. To make your 
    equations more compact, you may use the solidus (~/~), the exp function, or 
    appropriate exponents. Italicize Roman symbols for quantities and variables, 
    but not Greek symbols. Use a long dash rather than a hyphen for a minus 
    sign. Punctuate equations with commas or periods when they are part of a 
    sentence, as in:
    \begin{equation}
    a+b=\gamma\label{eq}
    \end{equation}
    
    Be sure that the symbols in your equation have been defined before or immediately following the equation. Use ``\eqref{eq}'', not ``Eq.~\eqref{eq}'' or ``equation \eqref{eq}'', except at 
    the beginning of a sentence: ``Equation \eqref{eq} is . . .''
    
    
    \section{\LaTeX-Specific Advice}
    
    Please use ``soft'' (e.g., \verb|\eqref{Eq}|) cross references instead
    of ``hard'' references (e.g., \verb|(1)|). That will make it possible
    to combine sections, add equations, or change the order of figures or
    citations without having to go through the file line by line.
    
    Please don't use the \verb|{eqnarray}| equation environment. Use
    \verb|{align}| or \verb|{IEEEeqnarray}| instead. The \verb|{eqnarray}|
    environment leaves unsightly spaces around relation symbols.
    
    Please note that the \verb|{subequations}| environment in {\LaTeX}
    will increment the main equation counter even when there are no
    equation numbers displayed. If you forget that, you might write an
    article in which the equation numbers skip from (17) to (20), causing
    the copy editors to wonder if you've discovered a new method of
    counting.
    
    {\BibTeX} does not work by magic. It doesn't get the bibliographic
    data from thin air but from .bib files. If you use {\BibTeX} to produce a
    bibliography you must send the .bib files. 
    
    {\LaTeX} can't read your mind. If you assign the same label to a
    subsubsection and a table, you might find that Table I has been cross
    referenced as Table IV-B3. 
    
    {\LaTeX} does not have precognitive abilities. If you put a
    \verb|\label| command before the command that updates the counter it's
    supposed to be using, the label will pick up the last counter to be
    cross referenced instead. In particular, a \verb|\label| command
    should not go before the caption of a figure or a table.
    
    Do not use \verb|\nonumber| inside the \verb|{array}| environment. It
    will not stop equation numbers inside \verb|{array}| (there won't be
    any anyway) and it might stop a wanted equation number in the
    surrounding equation.
    
    
    \section{Some Common Mistakes}\label{SCM}
    \begin{itemize}
    \item The word ``data'' is plural, not singular.
    \item The subscript for the permeability of vacuum $\mu_{0}$, and other common scientific constants, is zero with subscript formatting, not a lowercase letter ``o''.
    \item In American English, commas, semicolons, periods, question and exclamation marks are located within quotation marks only when a complete thought or name is cited, such as a title or full quotation. When quotation marks are used, instead of a bold or italic typeface, to highlight a word or phrase, punctuation should appear outside of the quotation marks. A parenthetical phrase or statement at the end of a sentence is punctuated outside of the closing parenthesis (like this). (A parenthetical sentence is punctuated within the parentheses.)
    \item A graph within a graph is an ``inset'', not an ``insert''. The word alternatively is preferred to the word ``alternately'' (unless you really mean something that alternates).
    \item Do not use the word ``essentially'' to mean ``approximately'' or ``effectively''.
    \item In your paper title, if the words ``that uses'' can accurately replace the word ``using'', capitalize the ``u''; if not, keep using lower-cased.
    \item Be aware of the different meanings of the homophones ``affect'' and ``effect'', ``complement'' and ``compliment'', ``discreet'' and ``discrete'', ``principal'' and ``principle''.
    \item Do not confuse ``imply'' and ``infer''.
    \item The prefix ``non'' is not a word; it should be joined to the word it modifies, usually without a hyphen.
    \item There is no period after the ``et'' in the Latin abbreviation ``et al.''.
    \item The abbreviation ``i.e.'' means ``that is'', and the abbreviation ``e.g.'' means ``for example''.
    \end{itemize}
    An excellent style manual for science writers is \cite{b7}.
    
    \section{Authors and Affiliations}
    \textbf{The class file is designed for, but not limited to, six authors.} A 
    minimum of one author is required for all conference articles. Author names 
    should be listed starting from left to right and then moving down to the 
    next line. This is the author sequence that will be used in future citations 
    and by indexing services. Names should not be listed in columns nor group by 
    affiliation. Please keep your affiliations as succinct as possible (for 
    example, do not differentiate among departments of the same organization).
    
    \section{Identify the Headings}
    Headings, or heads, are organizational devices that guide the reader through 
    your paper. There are two types: component heads and text heads.
    
    Component heads identify the different components of your paper and are not 
    topically subordinate to each other. Examples include Acknowledgments and 
    References and, for these, the correct style to use is ``Heading 5''. Use 
    ``figure caption'' for your Figure captions, and ``table head'' for your 
    table title. Run-in heads, such as ``Abstract'', will require you to apply a style (in this case, italic) in addition to the style provided by the drop down menu to differentiate the head from the text.
    
    Text heads organize the topics on a relational, hierarchical basis. For 
    example, the paper title is the primary text head because all subsequent 
    material relates and elaborates on this one topic. If there are two or more sub-topics, the next level head (uppercase Roman numerals) should be used and, conversely, if there are not at least two sub-topics, then no subheads should be introduced.
    
    
    \section{Figures and Tables}
    \paragraph{Positioning Figures and Tables} Place figures and tables at the top and bottom of columns. Avoid placing them in the middle of columns. Large figures and tables may span across both columns. Figure captions should be below the figures; table heads should appear above the tables. Insert figures and tables after they are cited in the text. Use the abbreviation ``Fig.~\ref{fig}'', even at the beginning of a sentence.
    
    \begin{table}[htbp]
    \caption{Table Type Styles}
    \begin{center}
    \begin{tabular}{|c|c|c|c|}
    \hline
    \textbf{Table}&\multicolumn{3}{|c|}{\textbf{Table Column Head}} \\
    \cline{2-4} 
    \textbf{Head} & \textbf{\textit{Table column subhead}}& \textbf{\textit{Subhead}}& \textbf{\textit{Subhead}} \\
    \hline
    copy& More table copy$^{\mathrm{a}}$& &  \\
    \hline
    \multicolumn{4}{l}{$^{\mathrm{a}}$Sample of a Table footnote.}
    \end{tabular}
    \label{tab1}
    \end{center}
    \end{table}
    
    
    \begin{figure}[htbp]
    \centerline{\includegraphics{fig1.png}}
    \caption{Example of a figure caption.}
    \label{fig}
    \end{figure}
    
    
    Figure Labels: Use words rather than symbols or abbreviations when writing Figure axis labels to avoid confusing the reader. As an example, write the quantity ``Magnetization'', or ``Magnetization, M'', not just ``M''. If including units in the label, present them within parentheses. Do not label axes only with units. In the example, write ``Magnetization (A/m)'' or ``Magnetization \{A[m(1)]\}'', not just ``A/m''. Do not label axes with a ratio of quantities and units. For example, write ``Temperature (K)'', not ``Temperature/K''.


    \pagestyle{plain}%
     %% 참고문헌
	\clearpage
	\phantomsection	
 
    \begin{thebibliography}{00}
    \bibitem{b1} G. Eason, B. Noble, and I. N. Sneddon, ``On certain integrals of Lipschitz-Hankel type involving products of Bessel functions,'' Phil. Trans. Roy. Soc. London, vol. A247, pp. 529--551, April 1955.
    \bibitem{b2} J. Clerk Maxwell, A Treatise on Electricity and Magnetism, 3rd ed., vol. 2. Oxford: Clarendon, 1892, pp.68--73.
    \bibitem{b3} I. S. Jacobs and C. P. Bean, ``Fine particles, thin films and exchange anisotropy,'' in Magnetism, vol. III, G. T. Rado and H. Suhl, Eds. New York: Academic, 1963, pp. 271--350.
    \bibitem{b4} K. Elissa, ``Title of paper if known,'' unpublished.
    \bibitem{b5} R. Nicole, ``Title of paper with only first word capitalized,'' J. Name Stand. Abbrev., in press.
    \bibitem{b6} Y. Yorozu, M. Hirano, K. Oka, and Y. Tagawa, ``Electron spectroscopy studies on magneto-optical media and plastic substrate interface,'' IEEE Transl. J. Magn. Japan, vol. 2, pp. 740--741, August 1987 [Digests 9th Annual Conf. Magnetics Japan, p. 301, 1982].
    \bibitem{b7} M. Young, The Technical Writer's Handbook. Mill Valley, CA: University Science, 1989.
    \end{thebibliography}

  % \appendix
  % \makeAppendix
  
  % \cedp
 
  %% 초록
  %% 국문초록

\begin{KoreanAbstract}
이 문서는 \LaTeX 에 대한 모델과 지침서입니다. 이와 seoultech.cls 파일은 귀하의 논문의 구성 요소들을 정의합니다 [제목, 본문, 제목, 등등]. *중요: 논문 제목이나 초록에서는 기호, 특수 문자, 각주, 또는 수학 기호를 사용하지 마십시오.
\end{KoreanAbstract}	
  \cedp
  
  % 감사의 글
  \markboth{감사의 글}{}
\begin{KoreanAcknowledgement}
.

\end{KoreanAcknowledgement}
  \cedp
  

\end{document}